\section{継続研鑽：活動・実績・証明をつなぐ}
\label{theory:profession:cpd}
CPDは専門職として能力を継続的に維持・向上させる活動である。仕事に従事している事実だけから、必要な学習と実績記録をすべて満たすとはしない。実務の中で何を学び、能力をどう改善したかを説明できる記録を残す\parencite{cpd_v3,cpd_versions,cpd_new2021}。

令和4年度問題には参加型、発信型、実務型、自己学習型という区分がある。研修の受講、技術指導や執筆、業務成果、自己学習は活動の性格が違う。換算条件と証明は制度・版ごとに確認する。学協会によって資格継続に確認条件があることと、国内の技術士資格が一律に自動失効することも区別する。

\section{2023コンピテンシーの八つのキーワード}
\label{theory:profession:competency2023}
2023年の資料では、専門的学識、問題解決、マネジメント、評価、コミュニケーション、リーダーシップ、技術者倫理、継続研さんの八つを挙げる。2014年の七つのキーワードから継続研さんを追加した時点を確認する\parencite{eth54_competency2023}。

知識を持つことだけでなく、複合的な問題を定義し、制約と利害を調べ、資源を配分し、結果を評価し、関係者と協力して責任ある成果を得る能力を考える。35歳程度での活躍への期待は、法定の受験年齢上限や資格の自動取消条件ではない。文書の版と発表主体を読み、2023版で独立しているコミュニケーションとリーダーシップを一つの項目にまとめない。
